% Generated draft: requires mathematical and source review.
\documentclass[11pt]{article}
\usepackage[margin=25mm]{geometry}
\usepackage{fontspec}
\setmainfont{Latin Modern Roman}
\usepackage{amsmath,amssymb,mathtools,unicode-math}
\setmathfont{Latin Modern Math}
\usepackage{enumitem,needspace,xurl,hyperref}
\setlength{\parindent}{0pt}
\setlength{\parskip}{5pt}
\setlength{\emergencystretch}{3em}
\newcommand{\N}{\mathbb{N}}
\newcommand{\Z}{\mathbb{Z}}
\newcommand{\Q}{\mathbb{Q}}
\newcommand{\R}{\mathbb{R}}
\newcommand{\C}{\mathbb{C}}
\newcommand{\F}{\mathbb{F}}
\newcommand{\A}{\mathbb{A}}
\newcommand{\PP}{\mathbb{P}}
\newcommand{\E}{\mathbb{E}}
\newcommand{\eps}{\varepsilon}
\newcommand{\dd}{\,\mathrm{d}}
\begin{document}
\section*{UnsolvedMath ID: 20001153}
\textbf{Source problem code:} catalogue entry
\subsection{Definitions and mathematical statement}
The original five-word fragment ``Conjecture 3 which in turn implies'' is not a standalone conjecture. The attached aid reports its source sentence as ``Conjecture 4 implies Conjecture 3 which in turn implies Conjecture 2.'' The full numbered conjectures are not included here; the aid extracts the following chromatic-to-divisibility implication from that context.

For a finite graph $F$ with at least one edge, let $\omega(F)$ be maximum clique size, $\chi(F)$ chromatic number, and $d(F)$ the least number of parts in a vertex partition such that no part contains a maximum clique of $F$. The aid states the exact rounded bound
\[
d(F)\le\left\lceil\frac{\chi(F)}{\omega(F)-1}\right\rceil.
\]
For triangle-free $F$ with an edge, $\omega(F)=2$ and avoiding a maximum clique means ordinary proper coloring, so equality holds: $d(F)=\chi(F)$.

A graph is $k$-divisible if every induced subgraph $H$ having an edge satisfies $d(H)\le k$. If a hereditary class satisfies $\chi(H)\le f(\omega(H))$, the reported transfer bound is
\[
k\le\sup_{\substack{w\ge2\\w\ \mathrm{attained\ as\ a\ clique\ number}}}
\left\lceil\frac{f(w)}{w-1}\right\rceil
\]
in the sense that a finite right side is a sufficient uniform divisibility constant. It need not be the smallest possible class constant. In particular $f(w)=2w-1$ gives $3$-divisibility. Maximum cliques here are distinguished from inclusionwise maximal cliques; the statement concerns ordinary divisibility, not strong clique-color divisibility.

The aid attributes to Chudnovsky--Seymour (2023) a bisimplicial-vertex theorem for even-hole-free graphs. Bisimplicial means the neighborhood is the union of two cliques; hereditary application yields the chromatic bound $\chi(H)\le2\omega(H)-1$ and therefore $3$-divisibility. Odd cycles of length at least five are even-hole-free and have $d=3$, so the constant three is sharp in that class. The aid reports that this makes the historical implication chain unconditional in its even-hole-free setting; exact source labels and bibliographic theorem details remain to be checked.
\subsection{Short English statement}
The record is a fragment of an implication chain. The supplied aid makes a rounded chromatic-to-divisibility transfer precise and applies a reported even-hole-free structural theorem, with odd cycles witnessing the sharp constant three.
\subsection{Sources}
\begin{itemize}
\item Local UnsolvedMath registry, permanent ID $20001153$, source code catalogue entry; canonical \texttt{aim-combinatorics-notes.json}, \texttt{notes[277]}.
\item AIM workshop \emph{The Perfect Graph Conjecture}: \url{https://aimath.org/WWN/perfectgraph/perfectgraph.pdf}. The canonical extraction is only the fragment quoted above; its complete sentence is supplied by the aid rather than independently inspected here. Original remarks and literature fields are empty.
\item Attached research attempt 1: context only, recorded run status not a problem, rigor label proved, novelty confidence low; registry status open, classification OPEN-TRIAGE, suggested difficulty L2. It names Chudnovsky and Seymour's 2023 bisimplicial-vertex theorem without a complete citation. The machine-generated aid is not peer reviewed; independently verify before citation.
\end{itemize}
\end{document}
